% Living Showcase: Stilmittel mit Nutzen (Palette Slot 8).
\StartChapter[ch:stilmittel]{Stilmittel}

\begin{runintext}
  Diese Mittel kosten keinen Platz, sparen aber Lesezeit. Sie lohnen sich
  überall dort, wo eine Umformung sonst Zeile für Zeile nachvollzogen
  werden müsste.
\end{runintext}

\SubsectionBar[sec:mhl]{Terme über Umformungen verfolgen}
\ZSFindex{Term-Einfärbung}

\begin{formulabox}[Wohin wandert welcher Koeffizient?]
  \[
    \ZSFmhlA{a}x^2 + \ZSFmhlB{b}x + \ZSFmhlD{c} = 0
  \]
  \ZSFsep
  \[
    \ZSFmhlC{x_{1,2}} = \frac{-\ZSFmhlB{b} \pm
      \sqrt{\ZSFmhlB{b}^2 - 4\,\ZSFmhlA{a}\,\ZSFmhlD{c}}}{2\,\ZSFmhlA{a}}
  \]
  \formulanote{Drei Gegebene, fünf Fundstellen im Resultat — mit Farbe ist
  die Zuordnung sofort da, ohne Farbe muss man sie jedes Mal nachlesen.}
\end{formulabox}

\begin{ZSFlist}
  \ZSFItem{ZSFmhlA/B/D}{Die gegebenen Grössen, je ein Strang.}
  \ZSFItem{ZSFmhlC}{Das Ziel — immer die Endform oder das Resultat.}
  \ZSFItem{Gleiche Farbe}{bedeutet immer dieselbe Grösse, nie blosse Betonung.}
\end{ZSFlist}

\SubsectionBar[sec:quantity]{Grössenfarben: Farbe als Identität}
\ZSFindex{Grössenfarbe}

\begin{runintext}
  Die Marker oben sind \ZSFkeyword*{positional} — sie gelten innerhalb einer
  Herleitung. Wo dieselben Grössen über das ganze Dokument hinweg auftauchen,
  gibt es das Gegenstück: eine Farbe gehört dauerhaft einer \ZSFkeyword{Grösse}.
\end{runintext}

\begin{formulabox}[Dieselbe Farbe, quer durch die ZSF]
  \[
    \ZSFqMoment{M} = \ZSFqKraft{F} \cdot \ZSFqWeg{r}
  \]
  \ZSFsep[Andere Stelle, andere Formel]
  \[
    W = \ZSFqKraft{F} \cdot \ZSFqWeg{s}
    \qquad
    \ZSFqKraft{F} = \frac{\ZSFqMoment{M}}{\ZSFqWeg{r}}
  \]
  \formulanote{Kraft, Weg und Moment behalten ihre Farbe in jeder Formel.
  Vergeben werden sie einmal in preamble.tex, nie im Kapitel.}
\end{formulabox}

\begin{defbox}[Zwei Systeme, zwei Fragen][weight=quiet]
  \ZSFkeyword*{ZSFmhlA..D} beantwortet »welche Rolle spielt dieser Term in
  \emph{dieser} Umformung«. \ZSFkeyword*{ZSFDeclareQuantity} beantwortet
  »welche Grösse ist das, überall«. Beide gleichzeitig zu verwenden ist
  möglich, aber selten sinnvoll — zwei Farbsprachen in einer Formel liest
  niemand unter Zeitdruck.
\end{defbox}

\SubsectionBar[sec:brace]{Formelteile benennen}
\ZSFindex{Klammer-Annotation}

\begin{formulabox}[Der Name steht am Teil, nicht daneben]
  \[
    \ZSFbraceunder{x^2+2x+1}{Ausgangsform}
    = \ZSFbraceover{(x+1)^2}{faktorisiert}
  \]
  \ZSFsep
  \[
    \ZSFsumAuto{k=1}{n} k
    = \ZSFbraceover{\frac{n(n+1)}{2}}{geschlossene Form}
  \]
  \formulanote{Eine Note unter der Box müsste erst sagen, welcher Teil
  gemeint ist. Die Klammer zeigt direkt darauf.}
\end{formulabox}

\SubsectionBar[sec:kombi]{Zusammenspiel}

\begin{formulabox}[Farbe, Klammer und Hervorhebung zusammen]
  \[
    \begin{aligned}
      \ZSFbraceunder{\ZSFmhlA{u}\,\ZSFmhlB{v}}{Produkt}\ '
      &= \ZSFmhlC{\ZSFmhlA{u}'\,\ZSFmhlB{v}} \\
      &\quad + \ZSFmhlC{\ZSFmhlA{u}\,\ZSFmhlB{v}'}
    \end{aligned}
  \]
  \ZSFhl{Farbe trägt die Zuordnung, die Klammer den Namen, die
  Hervorhebung die Kernaussage.}
\end{formulabox}

\begin{warnbox}[Wann es sich nicht lohnt]
  Farbe nur setzen, wenn die Zuordnung mathematisch eindeutig ist. Zur
  blossen Betonung eines Terms ist sie falsch — dafür gibt es
  \ZSFdanger{ZSFdanger} und ZSFhl. Im Zweifel keine Farbe.
\end{warnbox}
